В приложении представлен полный исходный код интеграционных тестов движка
нагрузочного тестирования из файла \texttt{tests/integration/attack\_test.go}.
Тесты используют библиотеку testcontainers-go для запуска реального
Docker-контейнера с сервисом httpbin и верифицируют корректность
метрик, классификации ошибок и обработки HTTP-статусов.

\begin{lstlisting}[language=Go,caption={Интеграционные тесты движка (\texttt{tests/integration/attack\_test.go})},label={lst:app-integration}]
//go:build integration

package integration_test

import (
    "testing"
    "time"

    "github.com/paniccaaa/tsunami/internal/attack"
)

func TestRunAttack_BasicSuccess(t *testing.T) {
    cfg := &attack.AttackConfig{
        URL:         testBaseURL + "/get",
        Method:      "GET",
        RPS:         10,
        Duration:    2 * time.Second,
        Workers:     5,
        Connections: 10,
        Timeout:     5 * time.Second,
    }
    stopCh := make(chan struct{})
    metrics, elapsed, err := attack.RunAttack(cfg, stopCh)
    if err != nil {
        t.Fatalf("RunAttack failed: %v", err)
    }
    if metrics.TotalRequests == 0 {
        t.Error("expected TotalRequests > 0")
    }
    if metrics.Failures > 0 {
        t.Errorf("expected 0 failures, got %d", metrics.Failures)
    }
    if metrics.StatusCodes[200] == 0 {
        t.Error("expected 200 status codes to be counted")
    }
    if metrics.TotalRequests != metrics.Successes {
        t.Errorf("all requests should be successful: total=%d successes=%d",
            metrics.TotalRequests, metrics.Successes)
    }
    if elapsed == 0 {
        t.Error("expected elapsed time > 0")
    }
    if metrics.MaxLatency == 0 {
        t.Error("expected MaxLatency to be recorded")
    }
    if metrics.MinLatency >= metrics.MaxLatency && metrics.TotalRequests > 1 {
        t.Errorf("expected MinLatency < MaxLatency, got min=%v max=%v",
            metrics.MinLatency, metrics.MaxLatency)
    }
}

func TestRunAttack_TimeoutClassification(t *testing.T) {
    cfg := &attack.AttackConfig{
        URL:         testBaseURL + "/delay/3",
        Method:      "GET",
        RPS:         5,
        Duration:    2 * time.Second,
        Workers:     5,
        Connections: 10,
        Timeout:     500 * time.Millisecond,
    }
    stopCh := make(chan struct{})
    metrics, _, err := attack.RunAttack(cfg, stopCh)
    if err != nil {
        t.Fatalf("RunAttack failed: %v", err)
    }
    if metrics.Failures == 0 {
        t.Error("expected failures due to timeout, got 0")
    }
    if metrics.Successes > 0 {
        t.Errorf("expected 0 successes, got %d", metrics.Successes)
    }
    if metrics.ErrorTypes[attack.ErrorTypeTimeout] == 0 {
        t.Errorf("expected ErrorTypeTimeout, got: %v", metrics.ErrorTypes)
    }
    for errType, count := range metrics.ErrorTypes {
        if errType != attack.ErrorTypeTimeout && count > 0 {
            t.Errorf("unexpected error type %q: %d", errType, count)
        }
    }
}

func TestRunAttack_MetricsAccuracy(t *testing.T) {
    t.Run("2xx is a success", func(t *testing.T) {
        cfg := &attack.AttackConfig{
            URL: testBaseURL + "/status/201", Method: "GET",
            RPS: 10, Duration: 2 * time.Second,
            Workers: 5, Connections: 10,
            Timeout: 5 * time.Second,
        }
        stopCh := make(chan struct{})
        metrics, _, err := attack.RunAttack(cfg, stopCh)
        if err != nil { t.Fatalf("RunAttack failed: %v", err) }
        if metrics.Successes == 0 {
            t.Error("expected Successes > 0 for /status/201")
        }
        if metrics.Failures > 0 {
            t.Errorf("expected 0 Failures for /status/201, got %d", metrics.Failures)
        }
        if metrics.StatusCodes[201] == 0 {
            t.Error("expected 201 to be tracked in StatusCodes")
        }
    })
    t.Run("5xx is a failure", func(t *testing.T) {
        cfg := &attack.AttackConfig{
            URL: testBaseURL + "/status/500", Method: "GET",
            RPS: 10, Duration: 2 * time.Second,
            Workers: 5, Connections: 10,
            Timeout: 5 * time.Second,
        }
        stopCh := make(chan struct{})
        metrics, _, err := attack.RunAttack(cfg, stopCh)
        if err != nil { t.Fatalf("RunAttack failed: %v", err) }
        if metrics.Failures == 0 {
            t.Error("expected Failures > 0 for /status/500")
        }
        if metrics.StatusCodes[500] == 0 {
            t.Error("expected 500 to be tracked in StatusCodes")
        }
        if len(metrics.ErrorTypes) > 0 {
            t.Errorf("HTTP 500 should not produce ErrorType entries, got: %v",
                metrics.ErrorTypes)
        }
    })
    t.Run("4xx is a failure", func(t *testing.T) {
        cfg := &attack.AttackConfig{
            URL: testBaseURL + "/status/404", Method: "GET",
            RPS: 10, Duration: 2 * time.Second,
            Workers: 5, Connections: 10,
            Timeout: 5 * time.Second,
        }
        stopCh := make(chan struct{})
        metrics, _, err := attack.RunAttack(cfg, stopCh)
        if err != nil { t.Fatalf("RunAttack failed: %v", err) }
        if metrics.Failures == 0 {
            t.Error("expected Failures > 0 for /status/404")
        }
        if metrics.StatusCodes[404] == 0 {
            t.Error("expected 404 to be tracked in StatusCodes")
        }
    })
}
\end{lstlisting}
